
\documentclass{article}

% ICLR 2027 style
\usepackage[final]{neurips_2024}  % Use latest available
\usepackage[utf8]{inputenc}
\usepackage[T1]{fontenc}
\usepackage{hyperref}
\usepackage{url}
\usepackage{booktabs}
\usepackage{amsfonts}
\usepackage{nicefrac}
\usepackage{microtype}
\usepackage{graphicx}
\usepackage{amsmath}
\usepackage{amssymb}
\usepackage{algorithm}
\usepackage{algorithmic}

\title{Quantum Arithmetic for Signal Classification: A PAC-Bayesian Framework with Geometric Interpretability}

\author{
  Anonymous Authors \\
  Paper under double-blind review for ICLR 2027
}

\begin{document}

\maketitle

\begin{abstract}
We introduce a novel signal classification framework based on Quantum Arithmetic (QA) - a modular arithmetic system with emergent geometric structure. Unlike black-box deep learning models, our approach provides: (1) geometric interpretability via algebraic topology, (2) PAC-Bayesian generalization guarantees, (3) 3000× parameter reduction vs CNNs, and (4) 10-100× inference speedup. We validate on seismic event classification and EEG seizure detection, achieving competitive accuracy with unique advantages in interpretability and efficiency.
\end{abstract}

\section{Introduction}
% Converted from markdown - TO BE FORMATTED

\section{Mathematical Framework}
% Converted from markdown - TO BE FORMATTED

\section{Seismic Event Classification}
% Converted from markdown - TO BE FORMATTED

\section{EEG Seizure Detection}
% Converted from markdown - TO BE FORMATTED

\section{Experimental Results}
% Converted from markdown - TO BE FORMATTED

\section{Discussion}
% Converted from markdown - TO BE FORMATTED

\section{Conclusion}
% Converted from markdown - TO BE FORMATTED

\bibliography{references}
\bibliographystyle{iclr2027}

\end{document}
